\subsection{序幺半群和序群的定义}

定义1. --- 集合M上的一个交换幺半群结构（写成加法形式）与一个序关系（写成 ≤ ）称为相容的，如果它们满足以下公理：

（OM）对每个 \(z \in \mathbf{M}\) ，关系 \(x \leqslant y\) 蕴含 \(x + z \leqslant y + z\) 。

配备了交换幺半群结构和与之相容的序关系的集合M称为序幺半群；若其交换幺半群结构是群结构，则称为序群。

用类似的方式可以定义非交换序幺半群的概念（VI，第30页，习题1）。

若一个序关系与给定的幺半群结构相容，则其相反序也与之相容。

例子. --- 1) 有理整数的加法群以及有理数的加法群，在I第21页与第117页所定义的序关系下都是序群。

* 实数的加法群亦然（一般拓扑，IV，第3页）。 *

\begin{enumerate}
\def\labelenumi{\arabic{enumi})}
\setcounter{enumi}{1}
\tightlist
\item
  * 定义在集合E上的有限数值函数的加法群，在序关系“ \(f(x) \leq g(x)\) 对所有 \(x \in E\) ”下是序群，该序关系记作 \(f \leq g\) 。这个关系表明函数 f 的图像位于函数 g 的图像之下。读者或许会发现，偶尔借助这一图像解释来思考是有益的。 *
\end{enumerate}

根据一般定义（集合论，IV，第264页），双射映射 \(f\) ——它把序幺半群 \(M\) 映到序幺半群 \(M'\) 上——称为 \(M\) 到 \(M'\) 上的一个同构，如果 \(M'\) 的结构是由 \(M\) 的结构沿 \(f\) 转移而得到的。这等价于说 \(f\) 是到 \(M'\) 上的一个映射，使得

\[
f (x + y) = f (x) + f (y)
\]

（也就是说，是幺半群 M 到幺半群 \(\mathbf{M}'\) 上的一个同态），并且关系 \(x \leqslant y\) 与 \(f(x) \leqslant f(y)\) 等价（由此特别地， \(f(x) = f(y)\) 蕴含 \(x = y\) ，也就是说 \(f\) 是单射）。

命题1（“不等式相加”）. --- 在序幺半群M中，设 \((x_{i})\) 与 \((y_{i})\) （ \(1 \leqslant i \leqslant n\) ）是两个由n个元素组成的序列，使得对所有 i 有 \(x_{i} \leqslant y_{i}\) ；则有

\[
x _ {1} + \dots + x _ {n} \leqslant y _ {1} + \dots + y _ {n}.
\]

若此外所有元素 \(x_{i}, y_{i}\) 都是可消去的（I，第15页，定义5）（特别地，若 \(M\) 是一个群），且 \(x_{i} < y_{i}\) 对某个 \(i\) 成立，则

\[
x _ {1} + \dots + x _ {n} <   y _ {1} + \dots + y _ {n}.
\]

一般情形通过归纳归结为n = 2的情形，利用可消去元素之和是可消去的这一事实来证明第二个断言（I，第15页，命题2）。第一个断言由以下关系推出

\[
x _ {1} + x _ {2} \leqslant x _ {1} + y _ {2} \quad \text { and } \quad x _ {1} + y _ {2} \leqslant y _ {1} + y _ {2},
\]

这些关系是假设以及(OM)的推论。在这种情况下，关系

\[
x _ {1} + x _ {2} = y _ {1} + y _ {2}
\]

将意味着

\[
x _ {1} + x _ {2} = x _ {1} + y _ {2} = y _ {1} + y _ {2},
\]

由此得出 \(x_{2} = y_{2}\) 与 \(x_{1} = y_{1}\) ，只要 \(x_{1}\) 与 \(y_{2}\) 是可消去的，这就证明了第二个断言。

命题2. --- 在序群G中，关系 \(x \leqslant y\) 与 \(x + z \leqslant y + z\) 等价。

事实上，通过在两边加上 z 或 \((-z)\) ，就可以由其中一个得到另一个。

这个事实可以表述为：在序群G中，序在平移下不变。换言之，平移是序群的序的一个自同构。

推论。 --- 在序群G中，关系 \(x \leqslant y\) 、\(0 \leqslant y - x\) 、\(x - y \leqslant 0\) 与 \(-y \leqslant -x\) 等价。

事实上，我们依次应用命题2，取 \(z = -x\) 、\(z = -y\) 与 \(z = - (x + y)\) 。

特别地，由这个推论可得，若G是序群，则从G到自身的映射 \(x \mapsto -x\) 把G的序变成相反的序。

